% CHAPTER 4 - IMPLEMENTATION AND TESTING
\chapter{Implementation and Testing}

\section{Tools and Technologies}
The implementation stack was chosen for clarity, type safety and the ability to run the
whole system on one development machine:

\begin{itemize}[leftmargin=*, topsep=4pt, itemsep=4pt]
    \item \textbf{FastAPI} \cite{fastapi}: request routing, dependency injection for database
          sessions, background tasks for memory extraction, automatic OpenAPI documentation,
          and a lifespan hook that checks the vector index at start-up.
    \item \textbf{SQLAlchemy 2} \cite{sqlalchemydocs}: declarative models for the platform
          database, and engines plus schema inspection for PostgreSQL, MySQL and SQLite target
          databases.
    \item \textbf{Pydantic v2 and pydantic-settings} \cite{pydanticdocs}: request and response
          schemas, validation of extracted memories, and typed configuration loaded from
          \texttt{.env}.
    \item \textbf{sqlglot} \cite{sqlglotdocs}: dialect-aware SQL parsing for the statement
          planner, and safe construction of preview and sampling queries.
    \item \textbf{sentence-transformers} \cite{sbertdocs}: the \texttt{all-MiniLM-L6-v2}
          embedding model and the \texttt{ms-marco-MiniLM-L-6-v2} cross-encoder, both loaded
          lazily on first use.
    \item \textbf{Qdrant} \cite{qdrantdocs} through \texttt{qdrant-client}: three collections
          with named dense and sparse vectors and server-side reciprocal rank fusion, embedded
          on disk at \texttt{./vector\_store} unless \texttt{VECTOR\_STORE\_URL} is set.
    \item \textbf{Google Gemini} \cite{geminidocs} through \texttt{google-generativeai}:
          separate model names for SQL generation, answer writing and memory extraction.
    \item \textbf{Jinja2 and htmx} \cite{htmxdocs}: server-rendered pages with partial
          updates; the only custom JavaScript (47 lines) handles panel toggling,
          Enter-to-send, copying SQL and closing modals.
    \item \textbf{pytest} \cite{pytestdocs} with FastAPI's test client for automated tests.
\end{itemize}

\section{Request Pipeline Implementation}
The \texttt{ask} handler in \texttt{app/api/chat.py} implements the flow of
Figure~\ref{fig:dataflow}. It validates that the question is not empty, authorizes the
user, resolves the target database and conversation, and stores the user message before
any LLM call so that the turn is never lost. Context building is wrapped in a
\texttt{try} block: if it fails, the handler logs a warning and calls a schema-only SQL
prompt instead, so a retrieval failure degrades quality but does not break the answer.

After SQL generation the handler calls \texttt{plan\_sql}. Refused statements return the
outcome \texttt{not\_allowed} with the planner's message; change statements are passed to
the change service; reads are executed by the connector. The answer is written by the
summary model, and if that call fails the handler falls back to \texttt{describe\_results},
which reports the row count and totals and averages of numeric columns. Finally the handler
saves the assistant message with its context trace, saves a \texttt{QueryMemory} whose SQL
template has string literals replaced by \texttt{'?'}, indexes it as an example, and
schedules memory extraction as a FastAPI background task.

Listing~\ref{lst:read} shows how reads are executed. The transaction is rolled back in a
\texttt{finally} block whatever the statement did, and one extra row is fetched to detect
truncation.

\begin{lstlisting}[language=Python, caption={Guarded read in \texttt{DatabaseConnector}}, label={lst:read}]
def read(self, engine: Engine, query: str, max_rows: int) -> tuple[list[dict], bool]:
    with engine.connect() as conn:
        transaction = conn.begin()
        try:
            result = conn.execute(text(query))
            if not result.returns_rows:
                return [], False
            columns = list(result.keys())
            fetched = result.fetchmany(max_rows + 1)
            rows = [dict(zip(columns, row)) for row in fetched[:max_rows]]
            return rows, len(fetched) > max_rows
        finally:
            transaction.rollback()
\end{lstlisting}

\subsection{Schema Introspection}
When a database is connected or first queried, the connector inspects every table's
columns and foreign keys and records each foreign key as a \texttt{references} entry such as
\texttt{customers.id}. It then samples distinct values of text columns (at most 200 tables,
12 columns per table, scanning 1,000 rows), skipping foreign-key and sensitive columns, and
keeps the list only if there are at most 20 distinct values of at most 40 characters each.
These values appear in the prompt (for example \texttt{status: delivered | pending |
shipped}), which helps the model write correct filter literals. The schema is cached under
the SHA-256 fingerprint of the connection identity and indexed into the
\texttt{schema\_tables} collection.

\section{Memory Store and Extraction}

\subsection{Memory Store}
\texttt{MemoryStore} implements creation, update, approval, rejection, expiry, superseding,
deletion, listing, statistics, version history and audit trail. Every operation first loads
the memory through an authorization check that returns HTTP 404 for unknown or deleted
memories and HTTP 403 for memories of another organization or another user's personal
memory. Creation checks for duplicates (identical normalized text or cosine similarity of
at least 0.96 with an active memory of the same user) and returns the existing memory
instead of creating a new one. Every state change is committed together with an audit
entry and then synchronized to the vector index, which keeps \texttt{approved},
\texttt{pending} and \texttt{superseded} memories and removes the others.

\subsection{Extraction}
\texttt{MemoryExtractor} sends the last user and assistant turn to the extraction model
with a prompt that defines the three memory types, forbids storing result rows, sensitive
values or credentials, forbids inferring rules only from generated SQL, and asks for strict
JSON. The response is parsed, invalid items are dropped, types, scopes and visibilities are
normalized to allowed values, and each item is validated with the \texttt{ExtractedMemory}
Pydantic model. Listing~\ref{lst:override} shows the override check that runs before the
LLM is called.

\begin{lstlisting}[language=Python, caption={One-time override check in \texttt{MemoryExtractor.extract}}, label={lst:override}]
latest_message = user_texts[-1]
if is_override(latest_message) and not is_rule_change(latest_message):
    logger.info("Detected one-query temporary override; skipping permanent rule extraction.")
    return []
\end{lstlisting}

During ingestion, an item marked as a correction with a subject triggers a conflict search:
approved or pending memories of the same user with the same subject whose embedding
similarity to the new content is below 0.85 are treated as contradicted. The first such
memory is superseded by the new one, which stays pending until reviewed.

\section{Hybrid Retrieval and the RAG Index}

\subsection{Vector Store}
\texttt{VectorStore} wraps the Qdrant client behind a lock and creates collections lazily
with a named 384-dimensional cosine dense vector and a named sparse vector with the IDF
modifier. Sparse vectors are built in \texttt{text.py}: each term is mapped to an index by
the first eight hex digits of its MD5 hash, and document term weights follow the BM25
saturation $tf\,(k_1+1)/(tf+k_1)$ with $k_1 = 1.2$, while query terms have weight 1. A
search issues one dense and one sparse prefetch with the same payload filter and fuses them
with \texttt{Fusion.RRF}. If the client cannot be created, the store reports itself
unavailable and every caller falls back to database scans.

\subsection{Index Synchronization}
\texttt{indexer.py} keeps three collections consistent with the platform database:
\texttt{memories} is updated on every memory change, \texttt{schema\_tables} receives one
point per table on schema sync (with a deterministic UUIDv5 identifier), and
\texttt{query\_examples} receives each successfully answered question with its masked SQL.
On start-up, \texttt{ensure\_index} compares each collection's point count with the
expected number of rows and rebuilds any collection that differs.

\subsection{Ranking}
The \texttt{MemoryRetriever.rank} method applies authorization filters, fetches vector
candidates, applies status and validity checks, removes duplicate content, computes the five
sub-scores of Section~3.5 and calls the reranker. Listing~\ref{lst:rerank} shows the
combination of the reranker score with the first-stage score.

\begin{lstlisting}[language=Python, caption={Combining the cross-encoder score with the first-stage score}, label={lst:rerank}]
weight = settings.RERANK_WEIGHT
for item, score in zip(scored, scores):
    item["first_stage_score"] = item["combined_score"]
    item["rerank_score"] = round(score, 4)
    item["combined_score"] = round(weight * score + (1 - weight) * item["combined_score"], 4)
\end{lstlisting}

Both the embedder and the reranker are wrapped in provider classes that can be replaced at
runtime. The test configuration replaces them with a deterministic 64-dimensional hash
embedder and a token-overlap scorer.

\section{Context Builder}
\texttt{ContextBuilder.build} returns the final prompt together with a
\texttt{context\_debug} record containing every retrieved memory with its five sub-scores
and rerank score, every excluded memory with its reason, the token budget and tokens used,
the selected schema tables and selection mode, the retrieval query actually used, the vector
store mode, and the examples included. This record is stored with the assistant message and
is what the Context Inspector displays. A second method builds the answer prompt: it
includes at most 20 result rows, per-column statistics (minimum, maximum, total and average
for numeric columns; distinct count and most common values otherwise), the business
definitions and preferences that were applied, and the list of masked columns, and asks for
one direct sentence followed by at most three insights without mentioning SQL.

\section{SQL Guard and Change Service}
\texttt{plan\_sql} returns a \texttt{SqlPlan} with the cleaned SQL, the kind
(\texttt{read}, \texttt{update}, \texttt{insert} or \texttt{delete}), the tables, the
target table and the parsed statement, or an error message. Listing~\ref{lst:plan} shows
the checks applied to updates and deletes.

\begin{lstlisting}[language=Python, caption={Checks for \texttt{UPDATE} and \texttt{DELETE} in \texttt{plan\_sql}}, label={lst:plan}]
if kind in ("update", "delete"):
    if statement.args.get("from") or statement.args.get("joins") or statement.args.get("using"):
        plan.error = "Changes that join other tables aren't supported. Name the rows directly."
        return plan
    if not statement.args.get("where"):
        plan.error = f"Refusing to {kind} every row in {plan.target}. Say which rows to change."
        return plan
\end{lstlisting}

\texttt{ChangeService.propose} refuses the change if data changes are disabled, if there is
no saved connection, or if the user lacks permission. It builds a \texttt{COUNT(*)} query and
a five-row sample query with the same \texttt{WHERE} condition using sqlglot (for inserts it
counts the listed \texttt{VALUES} or the rows of the source query), refuses zero-row and
over-limit changes, masks the sample for non-admins and stores the request. \texttt{approve}
re-plans the stored SQL, recounts, and calls \texttt{connector.write}, which raises
\texttt{RowCountMismatch} and rolls back if the driver reports a different number of
affected rows. Failures are stored on the request with status \texttt{failed}.

\section{Web Interface}
The web interface has a landing page where the visitor chooses to act as an Admin (may
request and apply changes) or an Employee (read-only; sensitive columns hidden), followed by
the database connection page, the SQL Assistant, the Memories dashboard, the Context
Inspector, the Knowledge page and the Evaluation page. State is held in three HTTP-only
cookies: the acting user, the saved connection and the last conversation. htmx requests
receive HTML fragments, and out-of-band swaps refresh the context panel, the conversation
list and the pending-memory badge. SQLite files can be uploaded (they must start with the
SQLite header and be at most 50\,MB) and are stored under \texttt{UPLOAD\_DIR} with a random
suffix. Figures~\ref{fig:ui_landing}--\ref{fig:ui_knowledge} are screenshots of the running
application.

\begin{figure}[H]
    \centering
    \figimg{0.48\textwidth}{ui_landing.png}{Landing page: choosing the Admin or Employee user type.}{fig:ui_landing}\hfill
    \figimg{0.48\textwidth}{ui_connect.png}{Database connection page (by path or upload).}{fig:ui_connect}
\end{figure}

\begin{figure}[H]
    \centering
    \figimg{0.48\textwidth}{ui_chat_change.png}{SQL Assistant workspace on the retail sample database.}{fig:ui_chat_empty}\hfill
    \figimg{0.48\textwidth}{ui_knowledge.png}{Knowledge page for entering a business rule manually.}{fig:ui_knowledge}
\end{figure}

The JSON API is documented automatically by FastAPI (Figure~\ref{fig:ui_swagger}); the full
endpoint list is given in Appendix~\ref{app:api}.

\begin{figure}[H]
    \centering
    \figimg{0.52\textwidth}{ui_swagger.png}{Automatically generated OpenAPI documentation of the JSON API.}{fig:ui_swagger}
\end{figure}

\section{Testing Methodology}

\subsection{Test Environment}
The automated tests in \texttt{tests/} run without network access, API keys or model
downloads. Fixtures in \texttt{conftest.py} replace the embedding model with a
64-dimensional hash embedder, the reranker with a token-overlap scorer and Qdrant with an
in-memory instance, and use an in-memory SQLite platform database. All LLM calls are
patched with fixed responses. The fixtures provide two organizations, users with different
roles and two database connections, which are used for isolation tests.

\subsection{Test Suite}
Table~\ref{tab:tests} summarises the suite. In total pytest collects 64 test cases in ten
files.

\begin{table}[H]
    \centering
    \small
    \caption{Automated Test Suite}
    \small
    \begin{tabular}{p{5.0cm} p{0.9cm} p{8.8cm}}
        \toprule
        \textbf{File} & \textbf{Tests} & \textbf{What is verified} \\
        \midrule
        \texttt{test\_pages.py}           & 19 & All pages render; upload accepts SQLite and rejects other files; chat renders answer and context; memory approval; change-request flow; first visit asks for user type; Employee is read-only. \\[3pt]
        \texttt{test\_changes.py}         & 12 & Planner classification; permission refusal; grant per table and action; approval required; rejected requests change nothing; approval stops if data changed; revoked grants; unsafe and oversized changes; masking and row cap; inserts and deletes. \\[3pt]
        \texttt{test\_memory\_store.py}   &  6 & Pending default; cross-org denial; versioning; approve/reject/delete; supersede; listing and statistics. \\[3pt]
        \texttt{test\_rag\_pipeline.py}   &  6 & Follow-up detection and rewritten retrieval; vector table search for large schemas; examples indexed and not shared between users; stale index rebuilt at start-up. \\[3pt]
        \texttt{test\_security.py}        &  6 & Cross-tenant, cross-user and connection isolation; viewer cannot share; SQL guard blocks writes; delete clears embedding. \\[3pt]
        \texttt{test\_failures.py}        &  4 & Extraction failure returns empty; chat survives retrieval failure; missing memory returns 404; code fences stripped. \\[3pt]
        \texttt{test\_retrieval.py}       &  4 & Only approved memories returned; inactive excluded; superseded recall; duplicates and conflicts. \\[3pt]
        \texttt{test\_extraction.py}      &  3 & One-time override not extracted; definition extracted; rule change detected. \\[3pt]
        \texttt{test\_api.py}             &  2 & End-to-end API demonstration flows; evaluation results endpoint. \\[3pt]
        \texttt{test\_context\_builder.py}&  2 & All prompt sections combined; token budget respected. \\
        \bottomrule
    \end{tabular}
    \label{tab:tests}
\end{table}

\subsection{Retrieval Benchmark}
In addition to the functional tests, \texttt{tests/run\_evaluation.py} builds a synthetic
dataset of 120 cases and measures retrieval and context assembly with the real MiniLM
embedder. Its design and results are presented in Chapter~5.
